\documentclass[12pt,a4paper]{article}

\usepackage[utf8]{inputenc}
\usepackage[T1]{fontenc}
\usepackage{amsmath, amssymb, amsthm}
\usepackage{geometry}
\usepackage{graphicx}
\usepackage{hyperref}
\usepackage{listings}
\usepackage{xcolor}
\usepackage{booktabs}
\usepackage{array}
\usepackage{caption}
\usepackage{enumitem}
\usepackage{bm}
\usepackage{mathtools}
\usepackage{algorithm}
\usepackage{algpseudocode}
\usepackage{float}

\geometry{margin=2.5cm}

\definecolor{codegray}{rgb}{0.95,0.95,0.95}
\definecolor{codeblue}{rgb}{0.0,0.3,0.6}
\definecolor{codegreen}{rgb}{0.1,0.5,0.1}
\definecolor{codepurple}{rgb}{0.4,0.0,0.6}

\lstset{
    backgroundcolor=\color{codegray},
    basicstyle=\ttfamily\small,
    keywordstyle=\color{codeblue}\bfseries,
    commentstyle=\color{codegreen}\itshape,
    stringstyle=\color{codepurple},
    breaklines=true,
    frame=single,
    framerule=0.5pt,
    rulecolor=\color{gray},
    numbers=left,
    numberstyle=\tiny\color{gray},
    numbersep=6pt,
    tabsize=4,
    showstringspaces=false,
    captionpos=b,
    language=Python
}

\hypersetup{
    colorlinks=true,
    linkcolor=codeblue,
    urlcolor=codeblue,
    citecolor=codeblue
}

\newcommand{\norm}[1]{\left\lVert#1\right\rVert}
\newcommand{\R}{\mathbb{R}}
\newcommand{\bx}{\mathbf{x}}
\newcommand{\bv}{\mathbf{v}}
\newcommand{\bg}{\mathbf{g}}
\newcommand{\bq}{\mathbf{q}}
\newcommand{\bz}{\mathbf{z}}
\newcommand{\bh}{\mathbf{h}}
\newcommand{\bp}{\mathbf{p}}
\newcommand{\bw}{\mathbf{w}}
\newcommand{\bomega}{\boldsymbol{\omega}}
\newcommand{\btheta}{\boldsymbol{\theta}}
\newcommand{\skew}[1]{\left\lfloor#1\right\rfloor_\times}

% ─────────────────────────────────────────────────────────────────────
\title{
    \Large\textbf{Pipeline 3 --- GPS-Free Drone Localisation}\\[0.5em]
    \large Complete Chronological Technical Breakdown\\[0.3em]
    \normalsize Error-State EKF $\cdot$ SuperPoint+LightGlue $\cdot$ Particle Filter $\cdot$
    Semantic Segmentation $\cdot$ Homography Localisation
}
\author{Emil Jagni\'c \\ \small All\_In\_One\_Pipeline / Pipeline\_3\_Rev1}
\date{April 2026}
% ─────────────────────────────────────────────────────────────────────

\begin{document}

\maketitle
\tableofcontents
\newpage

% ═══════════════════════════════════════════════════════════════════════════
\section{Architecture Overview}
% ═══════════════════════════════════════════════════════════════════════════

Pipeline 3 estimates drone GPS position without access to a GPS receiver. It
fuses three independent information sources:

\begin{enumerate}
    \item \textbf{Inertial Measurement Unit (IMU)} --- gyroscopes, accelerometers,
          barometer, airspeed, magnetometer, integrated via an Error-State Extended
          Kalman Filter (ES-EKF) for dead-reckoned position.
    \item \textbf{Visual feature matching} --- SuperPoint keypoints and LightGlue
          matches between the drone's camera frame and a precomputed TMS reference
          tile map, producing a sub-tile GPS measurement via homography.
    \item \textbf{Semantic terrain segmentation} --- a UNet++ semantic model
          classifies terrain classes in the query frame; class centroid alignment
          with precomputed tile predictions provides an independent confirmation
          of the feature-matching result.
\end{enumerate}

The end-to-end data flow for one frame is shown below.

\begin{figure}[H]
\centering
\begin{minipage}{0.82\textwidth}
\begin{verbatim}
t1: frame captured
  |
  +-- step_ekf(row, prev_ts)       --> ekf_state (lat, lon, yaw, sigma)
  |
  +-- rotate_image(frame, -yaw)    --> query_rotated  (north-aligned)
  |   resize to <= 1280 px         --> query_for_match
  |
  +-- preprocess(frame)            --> query_512x512
  |   semantic_model.predict()     --> class_mask (512x512 uint8)
  |
  +-- meta_tile_builder.run()
  |     find_tiles_within_radius() --> candidate tiles
  |     SP+LG match x N            --> first_pass ranked tiles
  |     8-neighbour expand top-1   --> second_pass tiles
  |     stitch top-3 --> meta_tile
  |     SP+LG verify(meta_tile)    --> verification_matches
  |
  +-- compute_dual_homography()    --> H, CShape, inliers
  |   extract_visual_measurements  --> homo_position
  |   look-ahead correction        --> homo_corrected
  |
  +-- semantic_confirmer.confirm() --> sem_conf
  |
  +-- if gate_pass:
  |     compute adaptive R
  |     ekf.update_position(homo_corrected, R)
  |
  +-- ekf.get_state()              --> (final_lat, final_lon, sigma)
t2 = t1 + dt: GPS estimate delivered  (dt ~ 2.8 s)
\end{verbatim}
\end{minipage}
\caption{Per-frame pipeline data flow (Pipeline 3).}
\end{figure}

\noindent The six pipeline phases are detailed in the following sections.

% ═══════════════════════════════════════════════════════════════════════════
\section{Phase 1 --- Error-State EKF (IMU Dead Reckoning)}
% ═══════════════════════════════════════════════════════════════════════════

\noindent\textbf{Primary file:} \texttt{src/ekf\_ins.py}

\subsection{Why an Error-State Formulation?}

A standard Extended Kalman Filter applied directly to a quaternion state requires
constrained optimisation ($\norm{\bq} = 1$). The error-state formulation avoids
this by separating the state into:
\begin{itemize}
    \item \textbf{Nominal state} $\tilde{\bq}$: the current best orientation
          estimate (quaternion), propagated through the nonlinear kinematics
          exactly.
    \item \textbf{Error state} $\delta\bx$: a small perturbation around the
          nominal state, maintained in a linear Kalman filter.
\end{itemize}

Because the error is small, the small-angle approximation
$\sin(\delta\theta) \approx \delta\theta$ holds, linearisation is accurate,
and the unit-norm constraint is trivially satisfied by normalising $\tilde{\bq}$
after each update. This follows Algorithm~4 from Kok, Hol, Sch\"{o}n
(2017)~\cite{kok2017}.

\subsection{State Vector}

The EKF maintains a $10$-dimensional error state:

\begin{equation}
    \delta\bx =
    \begin{bmatrix}
        \delta\btheta \\ \delta\mathbf{b}_g \\ \delta\bw \\ \delta\bp
    \end{bmatrix} \in \R^{10}
    \label{eq:error_state}
\end{equation}

where:
\begin{align*}
    \delta\btheta &\in \R^3 \quad \text{orientation error (small-angle, body frame)} \\
    \delta\mathbf{b}_g &\in \R^3 \quad \text{gyroscope bias error} \\
    \delta\bw &\in \R^2 \quad \text{wind velocity error (North--East, m/s)} \\
    \delta\bp &\in \R^2 \quad \text{horizontal position error (North--East, metres)}
\end{align*}

The corresponding $10 \times 10$ covariance matrix $P$ is initialised as:

\begin{equation}
    P_0 =
    \begin{bmatrix}
        0.5 \, I_3 & & & \\
        & 0.01 \, I_3 & & \\
        & & 25 \, I_2 & \\
        & & & 40000 \, I_2
    \end{bmatrix}
\end{equation}

The initial position uncertainty of $200$ m standard deviation ($200^2 = 40\,000$
m$^2$) reflects that the pipeline starts with no visual fix.

\subsection{MSFS Axis Mapping}

Microsoft Flight Simulator (MSFS) uses a \emph{left-handed} body frame:
$X_b = \text{right}$, $Y_b = \text{up}$, $Z_b = \text{forward}$.
Standard INS uses a \emph{right-handed} NED body frame:
$X_{NED} = \text{forward}$, $Y_{NED} = \text{right}$, $Z_{NED} = \text{down}$.

\paragraph{Accelerometer remapping.}
MSFS \texttt{ACCELERATION\_BODY\_*} variables return \emph{coordinate
acceleration} in ft/s$^2$ --- what an accelerometer with gravity blocked
would read. Real accelerometers output specific force
$\mathbf{f} = \mathbf{a}_{coord} + R_{nb}^\top \mathbf{g}_n$, where gravity
must be added back. Since MSFS provides accurate attitude angles, gravity is
synthesised directly:

\begin{equation}
    \mathbf{g}_b =
    \begin{bmatrix}
        -g \sin\theta \\ g \sin\phi\cos\theta \\ g\cos\phi\cos\theta
    \end{bmatrix}, \qquad
    \mathbf{f}_b =
    \underbrace{
        \begin{bmatrix} a_{Z,MSFS} \\ a_{X,MSFS} \\ -a_{Y,MSFS} \end{bmatrix}
    }_{\text{axis remap}}
    \cdot 0.3048 + \mathbf{g}_b
    \label{eq:accel_remap}
\end{equation}

where $\theta$ is pitch and $\phi$ is bank (roll), both in radians.
The factor $0.3048$ converts ft/s$^2$ to m/s$^2$.

\paragraph{Gyroscope remapping.}
Gyroscope readings are pseudovectors. Under the MSFS$\to$NED axis permutation,
the two sign changes (Y-negation for the right-hand rule, and handedness flip)
cancel exactly for pseudovectors. Therefore \emph{no sign negation} is applied:

\begin{equation}
    \bomega_{meas} =
    \begin{bmatrix} \omega_{Z,MSFS} \\ \omega_{X,MSFS} \\ \omega_{Y,MSFS} \end{bmatrix}
    \quad \text{(rad/s)}
\end{equation}

\subsection{Predict Step}

\begin{lstlisting}[caption={EKF predict step (ekf\_ins.py:213--247).}]
def predict(self, omega_meas, accel_body, dt):
    omega_corrected = omega_meas - self.gyro_bias
    dq = expq(dt/2 * omega_corrected)   # quaternion increment
    self.q_tilde = normalize(q_tilde x dq)
    R_nb = quat_to_rotation_matrix(self.q_tilde)

    F = eye(10)
    F[0:3, 3:6] = 0.5 * dt * R_nb    # orientation-bias coupling
    # position block: F[8:10, 8:10] = I_2  (already from eye)

    self.P = F @ P @ F.T + Q
    self.P[8:10, 8:10] += eye(2) * sigma_pos**2 * dt
\end{lstlisting}

\noindent The orientation is propagated via the quaternion exponential map:

\begin{equation}
    \exp_q(\bomega) =
    \begin{cases}
        \begin{bmatrix} \cos\frac{\theta}{2} \\[4pt]
        \dfrac{\sin(\theta/2)}{\theta}\,\bomega \end{bmatrix} & \theta > 0 \\[10pt]
        \begin{bmatrix} 1 \\ \mathbf{0} \end{bmatrix} & \theta = 0
    \end{cases}, \quad \theta = \norm{\bomega}
    \label{eq:expq}
\end{equation}

The state transition matrix $F$ propagates the covariance. The position error
block evolves as a random walk:

\begin{equation}
    P_{pos}^- = P_{pos} + \sigma_p^2 \, \Delta t \, I_2
    \label{eq:pos_rw}
\end{equation}

with $\sigma_p = 5.0$ m/$\sqrt{\text{s}}$. Uncertainty grows at approximately
$5\sqrt{\Delta t}$ metres per IMU step during dead reckoning.

Velocity is integrated in NED:

\begin{equation}
    \dot{\bv}_n = R_{nb}\,\mathbf{f}_b - \mathbf{g}_n
    \label{eq:vel_int}
\end{equation}

\subsection{Barometer Update}

The International Standard Atmosphere formula converts pressure to altitude:

\begin{equation}
    h = 44330 \left( 1 - \left(\frac{P}{P_0}\right)^{0.1903} \right), \quad
    P_0 = 1013.25 \text{ mbar}
    \label{eq:baro}
\end{equation}

Vertical velocity is estimated by finite differences:
$v_d = -(h_k - h_{k-1}) / \Delta t$. This update is performed outside the
Kalman framework --- altitude is treated as a direct measurement.

\subsection{Accelerometer and Magnetometer Update}

The accelerometer observation model exploits gravity as a reference vector.
In the body frame, the expected specific force (near-hover / slow manoeuvre)
is:

\begin{equation}
    \hat{\mathbf{f}}_b = R_{bn}\,\mathbf{g}_n
    \label{eq:accel_model}
\end{equation}

The linearised observation matrix for the orientation error states is:

\begin{equation}
    H_{accel} = R_{bn}\,\skew{\mathbf{g}_n}
    \label{eq:H_accel}
\end{equation}

where $\skew{\cdot}$ is the skew-symmetric matrix operator. The third column
of $H_{accel}$ is zeroed because yaw rotation about gravity is unobservable
from gravity direction alone.

During manoeuvres (detected by $\norm{a_h} > 2$ m/s$^2$ in the horizontal NED
plane), the accelerometer noise is inflated by factor 100 to de-weight the
update.

The magnetometer provides a scalar heading innovation:

\begin{equation}
    \tilde{y}_{mag} = \text{atan2}\!\left(\sin(\psi_{meas} - \hat{\psi}),\,
    \cos(\psi_{meas} - \hat{\psi})\right)
    \label{eq:heading_innov}
\end{equation}

Angle-wrapping ensures the innovation is always in $(-\pi, \pi]$. The
observation matrix selects only the yaw component from the error state:

\begin{equation}
    H_{mag} = \begin{bmatrix} 0 & 0 & 1 & 0 & 0 & 0 & 0 & 0 & 0 & 0 \end{bmatrix}
    \label{eq:H_mag}
\end{equation}

\subsection{Airspeed Update}

Airspeed $V_{air}$ is measured by the pitot tube. The observation model
decomposes it into NED components along the heading $\psi$:

\begin{equation}
    \bz =
    \begin{bmatrix} V_{air}\cos\psi \\ V_{air}\sin\psi \end{bmatrix} =
    \begin{bmatrix} v_n + w_n \\ v_e + w_e \end{bmatrix}
    \label{eq:airspeed_model}
\end{equation}

The Jacobian couples only wind states (indices 6, 7):

\begin{equation}
    H_{air} = \begin{bmatrix}
        0 & \cdots & 0 & 1 & 0 & 0 & 0 \\
        0 & \cdots & 0 & 0 & 1 & 0 & 0
    \end{bmatrix}_{2 \times 10}
\end{equation}

After the Kalman gain is applied, the velocity estimate is blended with the
airspeed-derived velocity using a factor $\alpha$ that adapts to flight
conditions:

\begin{equation}
    \alpha =
    \begin{cases}
        1.0 & \text{wind converged, no manoeuvre} \\
        0.8 & \text{wind not yet converged} \\
        0.5 & \text{active manoeuvre}
    \end{cases}
\end{equation}

\subsection{Visual Position Update}

This is the core closed-loop step that fuses the homography GPS measurement
into the EKF. The measurement $({\varphi}_{meas}, {\lambda}_{meas})$ in degrees is
converted to NED metres relative to the EKF reference origin
$({\varphi}_0, {\lambda}_0)$:

\begin{equation}
    \bz = \begin{bmatrix}
        (\varphi_{meas} - \varphi_0) \cdot R_E \\
        (\lambda_{meas} - \lambda_0) \cdot R_E \cos\varphi_0
    \end{bmatrix}
    \label{eq:visual_meas}
\end{equation}

The observation matrix selects the position error states at indices 8--9:

\begin{equation}
    H_{vis} =
    \begin{bmatrix}
        0 & \cdots & 0 & 1 & 0 \\
        0 & \cdots & 0 & 0 & 1
    \end{bmatrix}_{2 \times 10}
    \label{eq:H_vis}
\end{equation}

Standard Kalman measurement update:

\begin{align}
    \tilde{\mathbf{y}} &= \bz - \bh = \bz - \begin{bmatrix}p_n \\ p_e\end{bmatrix} \\
    S &= H_{vis} P H_{vis}^\top + R \\
    K &= P H_{vis}^\top S^{-1} \\
    \delta\bx &= K\,\tilde{\mathbf{y}} \\
    P^+ &= P - K S K^\top
    \label{eq:kalman}
\end{align}

\noindent The noise matrix $R = r_k I_2$ is set adaptively per frame (see
Section~\ref{sec:adaptive_R}).

% ═══════════════════════════════════════════════════════════════════════════
\section{Phase 2 --- Frame Capture and Preprocessing Bifurcation}
% ═══════════════════════════════════════════════════════════════════════════

\noindent\textbf{Primary files:} \texttt{runtime/run\_pipeline.py},
\texttt{src/image\_utils.py}

\subsection{Frame Acquisition}

\paragraph{SimConnect (live) mode.}
\texttt{SimConnectLiveSource.get\_latest\_frame()} returns
\texttt{(frame\_img, frame\_id, frame\_capture\_ts)}. The timestamp
$t_1 = $ \texttt{frame\_capture\_ts} is recorded immediately. The pipeline
returns a GPS estimate at $t_1 + \Delta t_{inf}$, during which the drone has
physically moved:

\begin{equation}
    d_{lag} = \int_{t_1}^{t_1 + \Delta t_{inf}} \norm{\bv(t)}\,dt
    \approx \norm{\bv} \cdot \Delta t_{inf}
    \label{eq:inference_lag}
\end{equation}

At $\Delta t_{inf} \approx 2.8$ s and $\norm{\bv} \approx 40$ m/s, this gives
$d_{lag} \approx 112$ m --- the camera look-ahead offset observed empirically
and corrected in Phase 6.

\paragraph{File mode.}
\texttt{load\_image(frame\_path)} reads a JPEG and returns a
\texttt{uint8} RGB array at original resolution. No resizing is performed at
this stage.

\subsection{Preprocessing Bifurcation}

The pipeline maintains two completely separate preprocessing paths to avoid
the ``black border'' keypoint problem:

\begin{table}[H]
\centering
\begin{tabular}{@{}lllll@{}}
\toprule
\textbf{Path} & \textbf{Input} & \textbf{Transform} & \textbf{Output} & \textbf{Consumer} \\
\midrule
Feature matching & 1920$\times$1079 RGB & Heading rotation + longest-edge resize & ${\leq}$1280$\times N$ RGB & SuperPoint+LightGlue \\
Semantic segm.  & 1920$\times$1079 RGB & Resize 512$\times$288, pad to 512$\times$512 & 512$\times$512 RGB & UNet++ \\
\bottomrule
\end{tabular}
\caption{The two preprocessing paths. Mixing them would cause black-border
keypoints (feature path) or incorrect class probabilities (semantic path).}
\end{table}

\paragraph{Why separate?}
SuperPoint was trained on natural images at their original resolution.
Padding introduces black borders that SuperPoint detects as high-contrast
edge keypoints --- garbage features at tile boundaries that poison the
homography estimate. The semantic UNet++ requires a fixed 512$\times$512
input by architecture design.

% ═══════════════════════════════════════════════════════════════════════════
\section{Phase 3a --- Semantic Terrain Segmentation}
% ═══════════════════════════════════════════════════════════════════════════

\noindent\textbf{Primary file:} \texttt{src/semantic\_model.py}

\subsection{Model Architecture}

The segmentation model is \textbf{UNet++} with an \textbf{EfficientNet-B3}
encoder, producing a 6-class pixel-wise segmentation mask.

\begin{table}[H]
\centering
\begin{tabular}{@{}rll@{}}
\toprule
\textbf{Class ID} & \textbf{Class} & \textbf{Colour (RGB)} \\
\midrule
0 & Waterbodies & (4, 4, 255) \\
1 & Forest / Trees & (0, 167, 2) \\
2 & Land & (243, 255, 150) \\
3 & Railway & (193, 105, 53) \\
4 & Roads & (255, 0, 231) \\
5 & Buildings & (150, 150, 150) \\
\bottomrule
\end{tabular}
\caption{Semantic terrain classes output by the UNet++ model.}
\end{table}

\subsection{Preprocessing for Semantic Path}

\begin{lstlisting}[caption={Semantic preprocessing (image\_utils.py).}]
def preprocess_query_frame(frame, resize_w=512, resize_h=288,
                            target_size=512):
    # Step 1: downsample preserving aspect ratio
    resized = cv2.resize(frame, (resize_w, resize_h))   # 512x288
    # Step 2: centre-pad to square
    pad_top = (target_size - resize_h) // 2  # = 112 px
    pad_bot = target_size - resize_h - pad_top
    padded = np.pad(resized, ((pad_top, pad_bot), (0,0), (0,0)))
    return padded  # 512x512
\end{lstlisting}

The vertical padding of $p_{top} = (512 - 288)/2 = 112$ px ensures the
drone's field-of-view is centred in the square input, preserving the
sky--ground geometry expected by the model's encoder.

\subsection{Semantic Histogram Pre-filter}

When \texttt{SEMANTIC\_PREFILTER\_ENABLED = True}, a
\texttt{SemanticTileScorer} computes histogram overlap between the query
semantic mask and each candidate tile's precomputed prediction mask
\emph{before} feature matching. The score:

\begin{equation}
    s_{hist} = \sum_{c=0}^{5} \min(h^{(q)}_c,\, h^{(r)}_c)
    \label{eq:hist_score}
\end{equation}

where $h^{(q)}_c$ and $h^{(r)}_c$ are normalised class-frequency histograms
for query and reference respectively. Only the top-$K = 10$ candidates by
$s_{hist}$ proceed to the expensive SuperPoint+LightGlue matching step.

% ═══════════════════════════════════════════════════════════════════════════
\section{Phase 3b --- Heading Rotation}
% ═══════════════════════════════════════════════════════════════════════════

\noindent\textbf{Primary file:} \texttt{src/visual\_measurement.py}

\subsection{Motivation}

Reference tiles are north-aligned orthophotos. A drone flying at heading
$\psi$ degrees sees the scene rotated by $\psi$ relative to the tile. Without
correction, SuperPoint keypoints in the query frame are spatially offset from
the corresponding reference keypoints by a rotation, which LightGlue cannot
reliably handle.

\subsection{Rotation with Canvas Expansion}

The query frame is rotated by $-\psi$ (i.e.\ the negative heading) to produce
a north-aligned view:

\begin{lstlisting}[caption={Heading rotation (visual\_measurement.py).}]
def rotate_image(img, angle_deg):
    h, w = img.shape[:2]
    cx, cy = w / 2, h / 2
    M = cv2.getRotationMatrix2D((cx, cy), angle_deg, 1.0)
    # Expand canvas to avoid cropping corners
    cos_a, sin_a = abs(M[0,0]), abs(M[0,1])
    new_w = int(h * sin_a + w * cos_a)
    new_h = int(h * cos_a + w * sin_a)
    M[0,2] += new_w/2 - cx
    M[1,2] += new_h/2 - cy
    rotated = cv2.warpAffine(img, M, (new_w, new_h))
    return rotated, M   # M stored for inverse transform
\end{lstlisting}

The canvas dimensions after rotation are:

\begin{equation}
    w_{new} = h|\sin\alpha| + w|\cos\alpha|, \quad
    h_{new} = h|\cos\alpha| + w|\sin\alpha|
    \label{eq:canvas_size}
\end{equation}

The rotation matrix $M$ is stored so that matched keypoints can later be
mapped back to the original frame coordinate system via the inverse transform,
if needed for the nadir correction.

\subsection{Performance Cap}

The rotated image is resized so its longest edge does not exceed
$d_{max} = 1280$ px:

\begin{equation}
    s = \frac{d_{max}}{\max(h_{rot}, w_{rot})}, \quad
    (w', h') = (\lfloor s\,w_{rot}\rfloor, \lfloor s\,h_{rot}\rfloor)
    \label{eq:resize}
\end{equation}

\texttt{INTER\_AREA} interpolation is used because it performs
area-weighted averaging during downscaling, preserving fine texture details
that SuperPoint relies on for repeatable keypoint detection. Bilinear
interpolation (\texttt{INTER\_LINEAR}) introduces aliasing artefacts on
downscaling.

% ═══════════════════════════════════════════════════════════════════════════
\section{Phase 4 --- Tile Search}
% ═══════════════════════════════════════════════════════════════════════════

\subsection{Frame 0: Best-First Search (Cold Start)}

\noindent\textbf{Primary file:} \texttt{src/best\_first\_search.py}

On the first frame, no particle filter exists. The
\texttt{BestFirstSearcher} exhaustively matches the rotated query against all
tiles within $r_{BFS} = 500$ m of the EKF position, using a priority queue
ordered by haversine distance.

\paragraph{Haversine distance.}
The great-circle distance between two GPS coordinates is:

\begin{equation}
    d = 2R_E \arcsin\!\sqrt{
        \sin^2\!\frac{\Delta\varphi}{2} +
        \cos\varphi_1\cos\varphi_2\,\sin^2\!\frac{\Delta\lambda}{2}
    }
    \label{eq:haversine}
\end{equation}

where $R_E = 6\,371\,000$ m. At TMS zoom 16 ($512$ px tiles), tile size is:

\begin{equation}
    T_{size} = \frac{2\pi R_E \cos\varphi}{2^{16}}
    \approx 332 \text{ m} \quad (\varphi \approx 55.6^\circ\text{N})
    \label{eq:tile_size}
\end{equation}

A 500 m radius covers approximately 7--12 tiles depending on exact position.

\subsection{Frame 1+: Particle-Guided Two-Pass Search}

\noindent\textbf{Primary files:} \texttt{src/particle\_filter.py},
\texttt{src/meta\_tile\_builder.py}

\subsubsection{Particle Filter Predict}

The particle filter maintains $N = 100$ particles
$\{(\bx_i, \hat{\psi}_i, w_i)\}_{i=1}^N$, where $\bx_i \in \R^2$ is position
in TMS tile coordinates. The motion model per particle:

\begin{align}
    \bx_i^- &= \bx_i +
        \frac{v \cdot \Delta t}{T_{size}}
        \begin{bmatrix} \sin\hat{\psi}_i \\ -\cos\hat{\psi}_i \end{bmatrix}
        + \mathcal{N}(0, \sigma_p^2)
    \label{eq:pf_predict_pos} \\
    \hat{\psi}_i^- &= \hat{\psi}_i + \dot{\psi}\,\Delta t
        + \mathcal{N}(0, \sigma_\psi^2)
    \label{eq:pf_predict_hdg}
\end{align}

where $v$ is speed (m/s), $\dot{\psi}$ is turn rate from the gyroscope, and
the $Y$-axis is negated because TMS Y increases southward (opposite to
geographic north). Process noise: $\sigma_p = 5$ m, $\sigma_\psi = 2^\circ$.

The search region radius is:

\begin{equation}
    r_{search} = \max\!\left( \sigma_{pos,PF} \cdot T_{size},\; r_{min} \right), \quad
    r_{min} = 500 \text{ m}
    \label{eq:search_radius}
\end{equation}

The $r_{min}$ floor guarantees that at least one full tile is always included
in the search, preventing the failure mode where the EKF position lands near a
tile boundary and no tiles are returned.

\subsubsection{First-Pass Search}

\texttt{MetaTileBuilder.run()} calls \texttt{find\_tiles\_within\_radius()}
to enumerate all tiles within $r_{search}$ of the EKF position, then runs
SuperPoint+LightGlue on each. The result is a sorted list:

\begin{equation}
    \mathcal{T}_{1st} = \{(tx_k, ty_k, n_k)\}_{k=1}^{M}
    \quad \text{sorted by } n_k \text{ (match count) descending}
    \label{eq:first_pass}
\end{equation}

\subsubsection{Second-Pass: 8-Neighbour Expansion}

The top-1 tile from the first pass is expanded to its 8-connected neighbours:

\begin{equation}
    \mathcal{T}_{2nd} = \{(tx^* + \Delta x,\; ty^* + \Delta y) :
    (\Delta x, \Delta y) \in \{-1,0,1\}^2 \setminus \{(0,0)\}\}
    \label{eq:second_pass}
\end{equation}

These 8 tiles are also matched against the query. This ensures that if the
ground truth straddles a tile boundary, the correct tile is captured in the
candidate set.

\subsubsection{Meta-Tile Construction and Verification}

The top-3 tiles by match count are stitched into a single reference image
(the \emph{meta-tile}). The full query frame is then re-matched against this
meta-tile in a \emph{verification match}. A meta-tile is flagged as
\texttt{verified} if:

\begin{equation}
    n_{verify} \geq T_{match} = 25
    \label{eq:verify_threshold}
\end{equation}

\paragraph{Why meta-tiles?}
A single $512 \times 512$ tile covers $\approx 332 \times 332$ m. The drone's
field of view at 1500 ft altitude is wider than one tile. By stitching 3 tiles,
the reference image provides a $\approx 1000 \times 332$ m footprint, more
match points, and a larger spatial baseline for the homography, reducing
sensitivity to localised repetitive textures.

\subsection{TMS Coordinate System}

TMS (Tile Map Service) tiles use an inverted Y-axis compared to OSM:

\begin{equation}
    y_{TMS} = 2^z - 1 - y_{OSM}
    \label{eq:tms_y}
\end{equation}

This inversion was the source of Bug~1 (see CLAUDE.md). All tile-coordinate
functions in \texttt{tile\_utils.py} apply this conversion consistently.

% ═══════════════════════════════════════════════════════════════════════════
\section{Phase 5 --- Feature Matching (SuperPoint + LightGlue)}
% ═══════════════════════════════════════════════════════════════════════════

\noindent\textbf{Primary file:} \texttt{src/geometric\_matcher.py}

\subsection{SuperPoint Keypoint Detector and Descriptor}

SuperPoint~\cite{detone2018} is a self-supervised CNN trained to detect and
describe keypoints simultaneously. It outputs:

\begin{itemize}
    \item \textbf{Keypoint heatmap} $H \in [0,1]^{H/8 \times W/8}$:
          probability of a keypoint at each cell
    \item \textbf{Descriptor map} $D \in \R^{256 \times H \times W}$:
          dense 256-dimensional descriptor at every pixel
\end{itemize}

Keypoints are extracted as local heatmap maxima above a confidence threshold,
then non-maximum suppressed to at most \texttt{MAX\_NUM\_KEYPOINTS} $= 2048$.

\paragraph{Why SuperPoint over SIFT/ORB?}
Traditional detectors (SIFT, ORB) rely on handcrafted gradient histograms that
are sensitive to the large illumination and style gap between MSFS 3D rendering
and orthophoto satellite imagery. SuperPoint's learned descriptors generalise
better to this domain shift because they were trained on a diverse augmentation
pipeline including simulated homographic warps and photometric noise.

\subsection{LightGlue Matcher}

LightGlue~\cite{lindenberger2023} is a transformer-based matcher that computes
attention-weighted feature correspondences end-to-end. Given two sets of
SuperPoint keypoints and descriptors, it outputs:

\begin{itemize}
    \item \texttt{matches}: $K \times 2$ array of index pairs $(i, j)$
    \item \texttt{matching\_scores}: $K$-vector of match confidence values
\end{itemize}

Unlike descriptor-distance matchers (nearest-neighbour + ratio test),
LightGlue uses cross-attention to reason globally about the geometric
consistency of the match set. This makes it significantly more robust
under partial overlap and appearance change.

% ═══════════════════════════════════════════════════════════════════════════
\section{Phase 6 --- Homography and GPS Estimation}
% ═══════════════════════════════════════════════════════════════════════════

\noindent\textbf{Primary file:} \texttt{src/visual\_measurement.py}

\subsection{Dual Homography Estimation}

Two independent RANSAC-based homography estimates are computed on the same
match set:

\paragraph{MAGSAC++~\cite{barath2020}.}
A probabilistic RANSAC variant that weights each sample by its probability of
being an inlier under the model, avoiding a hard inlier threshold. Tends to be
more accurate when inlier ratios are moderate ($30$--$70\%$).

\paragraph{LMEDS.}
Least-Median-of-Squares estimator. Minimises the median squared reprojection
error. Computationally faster, but sensitive to high outlier ratios ($> 50\%$).

The homography $H$ maps query-frame pixel coordinates to reference-tile pixel
coordinates:

\begin{equation}
    \begin{bmatrix}u_r \\ v_r \\ 1\end{bmatrix} \sim
    H \begin{bmatrix}u_q \\ v_q \\ 1\end{bmatrix}
    \label{eq:homography}
\end{equation}

\subsection{CShape Quality Score}

The \emph{CShape} score measures how well the homography maps the query image
corners to a valid convex quadrilateral in the reference:

\begin{equation}
    \text{CShape} =
    \frac{\text{area}\!\left(\text{ConvexHull}(H\,\mathbf{c}_1, \ldots, H\,\mathbf{c}_4)\right)}
    {\text{area}\!\left(\text{BoundingBox}(H\,\mathbf{c}_1, \ldots, H\,\mathbf{c}_4)\right)}
    \label{eq:cshape}
\end{equation}

where $\mathbf{c}_1,\ldots,\mathbf{c}_4$ are the four query-frame corners.
A perfect planar homography maps a rectangle to a convex quadrilateral
($\text{CShape} = 1$). A degenerate or folded homography produces a
non-convex result ($\text{CShape} < 0.3$), which is rejected by the quality
gate.

The winner between MAGSAC++ and LMEDS is selected as:

\begin{equation}
    k^* = \arg\max_{k \in \{\text{MAGSAC, LMEDS}\}}
    \left( \text{CShape}_k \times n_{inliers,k} \right)
\end{equation}

\subsection{Visual Measurement Cascade}

Five methods extract a GPS position from the winning homography $H$:

\begin{table}[H]
\centering
\begin{tabular}{@{}rll@{}}
\toprule
\textbf{Priority} & \textbf{Method} & \textbf{Description} \\
\midrule
1 & \texttt{nadir\_corrected} & Project nadir point corrected for pitch and roll \\
2 & \texttt{trimmed\_centroid} & Centroid of inliers after 10\% trimming \\
3 & \texttt{inlier\_centroid} & Mean of all inlier reference positions \\
4 & \texttt{weighted\_centroid} & Inliers weighted by match confidence score \\
5 & \texttt{projected\_center} & Map frame centre through $H$ \\
\bottomrule
\end{tabular}
\caption{Visual measurement cascade --- first valid result is used.}
\end{table}

\paragraph{Nadir-corrected projection.}
The drone's camera has a fixed forward tilt. Pitch $\theta$ and roll $\phi$
displace the ground-point directly below the drone from the image centre.
The nadir pixel in the query frame is:

\begin{equation}
    \mathbf{p}_{nadir} =
    \begin{bmatrix} c_x - f\tan\phi \\ c_y + f\tan\theta \end{bmatrix}
    \label{eq:nadir}
\end{equation}

Applying $H$ to $\mathbf{p}_{nadir}$ gives the reference-tile pixel
corresponding to the true ground position below the drone.

\subsection{Tile Pixel to GPS Conversion}

Given reference-tile pixel $(u_r, v_r)$ in a tile $(tx, ty)$ at TMS zoom
$z$, the GPS coordinates are:

\begin{align}
    \lambda &= \lambda_{SW}(tx,ty,z) + \frac{u_r}{512}\,\Delta\lambda_{tile}
    \label{eq:pix2lon} \\
    \varphi &= \varphi_{NE}(tx,ty,z) - \frac{v_r}{512}\,\Delta\varphi_{tile}
    \label{eq:pix2lat}
\end{align}

The sign convention on $v_r$ reflects the TMS Y-down screen coordinate system
($v_r = 0$ is the tile's north edge).

\subsection{Camera Look-Ahead Correction}

Empirical diagnostic analysis revealed that all homography positions were
systematically $L = 110$ m \emph{ahead} of the drone's true position in the
heading direction. The MSFS virtual camera has a fixed forward tilt, so the
image centre corresponds to ground ahead of the drone rather than directly
below it (the nadir correction alone is insufficient because the tilt is
baked into the camera model, not exposed as a pitch angle).

The correction shifts the position backward along the heading direction:

\begin{align}
    L_{eff} &= L \cdot \cos\phi
    \label{eq:lookahead_eff} \\
    \Delta N &= -L_{eff}\cos\psi, \quad \Delta E = -L_{eff}\sin\psi
    \label{eq:lookahead_ned} \\
    \varphi_{corr} &= \varphi_{homo} + \frac{\Delta N}{111320}
    \label{eq:lookahead_lat} \\
    \lambda_{corr} &= \lambda_{homo} + \frac{\Delta E}{111320\cos\varphi_{homo}}
    \label{eq:lookahead_lon}
\end{align}

The $\cos\phi$ factor (bank angle) reduces the effective look-ahead during
banked turns because the camera's forward projection component decreases as
the aircraft rolls. This correction reduced mean position error from 103.5 m
to 9.7 m (91\% improvement).

% ═══════════════════════════════════════════════════════════════════════════
\section{Phase 7 --- Semantic Double-Confirmation}
% ═══════════════════════════════════════════════════════════════════════════

\noindent\textbf{Primary file:} \texttt{src/semantic\_confirmer.py}

The semantic confirmer provides an independent verification of the tile match
that is completely decoupled from feature matching. For each semantic class $c$
that appears significantly in both the query mask and the reference tile
prediction mask, the spatial centroids are computed and compared:

\begin{equation}
    d_c = \norm{\bar{\mathbf{p}}^{(q)}_c - \bar{\mathbf{p}}^{(r)}_c}_2
    \label{eq:centroid_dist}
\end{equation}

A class pair is declared aligned if $d_c < 50$ px. The confirmation confidence
is the aligned fraction:

\begin{equation}
    c_{sem} =
    \frac{\left|\left\{c : d_c < 50 \text{ px}\right\}\right|}
    {|\text{significant classes}|}
    \label{eq:sem_conf}
\end{equation}

\paragraph{Purpose.}
A high CShape + high inlier count can still arise from a \emph{wrong} tile if
the scene contains repetitive patterns (e.g.\ a regular grid of buildings).
The semantic confirmation catches this case: even if feature matching is
geometrically consistent, the terrain class layout provides an independent
spatial vote.

% ═══════════════════════════════════════════════════════════════════════════
\section{Phase 8 --- Adaptive Measurement Noise and Final Estimate}
% ═══════════════════════════════════════════════════════════════════════════
\label{sec:adaptive_R}

\noindent\textbf{Primary file:} \texttt{runtime/run\_pipeline.py:376--390}

The measurement noise variance $r_k$ is computed adaptively from four
independent quality signals:

\begin{equation}
    r_k = r_{base} \cdot m_{bank} \cdot m_{verify} \cdot m_{sem}
    \label{eq:adaptive_R}
\end{equation}

where:

\begin{align}
    r_{base} &=
    \begin{cases}
        30^2 = 900 \text{ m}^2 & \text{CS} > 0.5 \text{ and } n_{inl} > 100 \\
        60^2 = 3600 \text{ m}^2 & \text{otherwise} \\
        100^2 = 10000 \text{ m}^2 & \text{cold start}
    \end{cases} \\[6pt]
    m_{bank} &=
    \begin{cases}
        2.0 & |\phi| > 0.35 \text{ rad} \; (\approx 20^\circ) \\
        1.0 & \text{otherwise}
    \end{cases} \\[6pt]
    m_{verify} &=
    \begin{cases}
        2.0 & \text{meta-tile not verified} \\
        1.0 & \text{meta-tile verified}
    \end{cases} \\[6pt]
    m_{sem} &= \max(0.5,\; 2.0 - 1.5\,c_{sem})
\end{align}

The semantic factor $m_{sem}$ decreases $r_k$ (increases EKF trust) as
$c_{sem}$ approaches 1.0, and increases $r_k$ (decreases trust) when
$c_{sem}$ is low:

\begin{equation}
    c_{sem} = 1.0 \;\Rightarrow\; m_{sem} = 0.5 \quad \text{(half noise)}
    \qquad
    c_{sem} = 0.0 \;\Rightarrow\; m_{sem} = 2.0 \quad \text{(double noise)}
\end{equation}

\subsection{Quality Gate}

A frame triggers an EKF update only if all three conditions are met:

\begin{equation}
    \text{gate\_pass} = \mathbf{1}\!\left[
        \text{CShape} > 0.3 \;\wedge\;
        n_{inliers} > 20 \;\wedge\;
        \text{homo\_position} \neq \emptyset
    \right]
    \label{eq:quality_gate}
\end{equation}

Frames that fail the gate cause the EKF to coast on IMU dead reckoning, growing
$\sigma_{pos}$ until the next successful gate pass.

\subsection{Final State Extraction}

The final GPS estimate is extracted from the EKF:

\begin{align}
    \varphi_{final} &= \varphi_0 + \frac{p_n}{R_E}
    \label{eq:final_lat} \\
    \lambda_{final} &= \lambda_0 + \frac{p_e}{R_E\cos\varphi_0}
    \label{eq:final_lon} \\
    \sigma_{pos} &= \sqrt{\max(P[8,8],\, P[9,9])}
    \label{eq:sigma_pos}
\end{align}

The $\sigma_{pos}$ value is the EKF's self-assessed 1-sigma position
uncertainty in metres. In steady-state operation it oscillates between
$\approx 10$ m (just after a high-quality gate pass) and $\approx 60$ m
(several seconds of dead reckoning at $\sigma_p = 5$ m/$\sqrt{\text{s}}$).

% ═══════════════════════════════════════════════════════════════════════════
\section{Particle Filter Update and Resample}
% ═══════════════════════════════════════════════════════════════════════════

\noindent\textbf{Primary file:} \texttt{src/particle\_filter.py}

The particle filter is used \emph{solely} to determine the search region for
the next frame. The EKF is the authoritative position estimator.

\subsection{Weight Update}

Each particle weight is multiplied by the likelihood of the observed tile
match position:

\begin{equation}
    w_i^+ \propto w_i \cdot
    \exp\!\left(-\frac{\norm{\bx_i - \bx_{meas}}^2}{2\sigma_m^2}\right)
    \label{eq:pf_update}
\end{equation}

where $\sigma_m = \texttt{MEASUREMENT\_NOISE\_POSITION\_M} = 500$ m. The large
measurement noise reflects that domain shift makes tile matches unreliable;
the EKF (with its adaptive $R$) is trusted for position, while the particle
filter only provides a coarse search window.

\subsection{Effective Sample Size and Resampling}

The effective sample size:

\begin{equation}
    N_{eff} = \frac{1}{\sum_i w_i^2}
    \label{eq:neff}
\end{equation}

When $N_{eff} < \texttt{RESAMPLE\_THRESHOLD} \times N = 0.5 \times 100 = 50$,
systematic resampling is triggered to prevent weight degeneracy. Particle
spread $\sigma_{PF}$ is computed as the standard deviation of particle
positions in metres, reported per frame for diagnostics.

\subsection{Divergence Detection}

If the particle spread exceeds $\texttt{DIVERGENCE\_POSITION\_THRESHOLD} = 500$ m
or the maximum particle weight falls below $\texttt{DIVERGENCE\_WEIGHT\_THRESHOLD} = 0.01$,
divergence is flagged. On the next call to \texttt{process\_frame},
\texttt{frame\_count} is reset to~0, triggering a cold-start BFS reinitialisation.

% ═══════════════════════════════════════════════════════════════════════════
\section{Output Files and Data Flags}
% ═══════════════════════════════════════════════════════════════════════════

\noindent\textbf{Primary file:} \texttt{config/config.py}

\begin{table}[H]
\centering
\begin{tabular}{@{}lll@{}}
\toprule
\textbf{Flag} & \textbf{Default} & \textbf{Output} \\
\midrule
\texttt{SAVE\_QUERY\_FRAMES} & False & \texttt{flight\_data/frame\_NNNN.jpg} \\
\texttt{SAVE\_IMU\_ROWS} & False & \texttt{flight\_data/frame\_NNNN\_imu.json} \\
\texttt{SAVE\_ANALYSIS\_DATA} & False & \texttt{px4\_gps\_input.csv}, \texttt{analysis\_extras.csv} \\
\texttt{SAVE\_TIMING\_DATA} & False & \texttt{timing\_data.csv} \\
\texttt{SAVE\_PIPELINE\_TRACE} & True & \texttt{pipeline\_data/frame\_NNNN/} (7 files) \\
\bottomrule
\end{tabular}
\caption{Config flags controlling per-frame data capture.}
\end{table}

\paragraph{PX4 GPS\_INPUT (MAVLink MSG~232).}
When \texttt{SAVE\_ANALYSIS\_DATA = True}, the pipeline writes a
\texttt{px4\_gps\_input.csv} compatible with PX4's GPS\_INPUT MAVLink
message~(ID~232). Key field mappings:

\begin{equation*}
    \texttt{lat} = \lfloor\varphi_{final} \times 10^7\rfloor,\quad
    \texttt{lon} = \lfloor\lambda_{final} \times 10^7\rfloor,\quad
    \texttt{yaw} = \lfloor\psi \times 100\rfloor
\end{equation*}

Fields \texttt{hdop} and \texttt{vdop} are set to~0 and masked via
\texttt{ignore\_flags = 0x0006} (bit~1 = ignore hdop, bit~2 = ignore vdop).

% ═══════════════════════════════════════════════════════════════════════════
\section{Experimental Results}
% ═══════════════════════════════════════════════════════════════════════════

Results on run \texttt{live\_015\_cph\_1500ft} (63 frames, CPH reference map,
1500 ft altitude):

\begin{table}[H]
\centering
\begin{tabular}{@{}lrrrrrrr@{}}
\toprule
\textbf{Method} & $n$ & \textbf{Mean} & \textbf{Med.} & \textbf{Std.} & \textbf{Max}
    & $<$50\,m & $<$100\,m \\
\midrule
Online EKF (visual+IMU) & 63 & 35.2\,m & 32.9\,m & 22.8\,m & 87\,m & 71.4\% & 100\% \\
Homo raw (uncorrected) & 63 & 167.7\,m & 116.3\,m & 120.3\,m & 622\,m & 0.0\% & 20.6\% \\
Homo corrected (+110\,m) & 63 & 86.4\,m & 20.8\,m & 138.9\,m & 701\,m & 69.8\% & 73.0\% \\
\bottomrule
\end{tabular}
\caption{Multi-method accuracy comparison. The online EKF combines the
look-ahead-corrected homography with IMU dead reckoning via the 10D
error-state filter, achieving the best mean and worst-case error.}
\end{table}

Gate pass rate: 57/63 (90\%). Method distribution: 1 cold start, 62 temporal
tracking. All 63 frames were within the reference map.

% ═══════════════════════════════════════════════════════════════════════════
\section{Known Limitations and Future Work}
% ═══════════════════════════════════════════════════════════════════════════

\begin{enumerate}
    \item \textbf{Camera look-ahead constant.}
        $L = 110$ m is empirically calibrated for this specific camera model
        and altitude. A proper camera calibration (intrinsics + extrinsics)
        would make it altitude- and tilt-independent.

    \item \textbf{Domain shift.}
        MSFS renders 3D geometry with dynamic lighting, while reference tiles
        are flat 2D orthophotos. SuperPoint descriptors generalise, but
        CShape degrades near terrain transitions (forest edges, coastlines).

    \item \textbf{Speed.}
        2.8\,s/frame ($\approx 0.36$ Hz) vs.\ a typical camera rate of 30 Hz.
        Bottleneck: SuperPoint inference on the rotated frame + LightGlue.
        Potential improvement: cache the query SuperPoint features across all
        tile matches in one pass (currently re-run per tile pair).

    \item \textbf{Repeated-structure ambiguity.}
        Dense urban grids can produce geometrically consistent but positionally
        wrong homographies. Semantic confirmation mitigates this but does not
        fully eliminate it.

    \item \textbf{Out-of-map handling.}
        When the drone exits the reference tile map, the cold-start BFS finds
        0 tiles and the system falls back to IMU dead reckoning indefinitely.
        A map-extent check on the EKF position would give earlier warning.
\end{enumerate}

% ═══════════════════════════════════════════════════════════════════════════
\begin{thebibliography}{9}

\bibitem{kok2017}
M.~Kok, J.~D.~Hol, and T.~B.~Sch\"{o}n,
``Using Inertial Sensors for Position and Orientation Estimation,''
\textit{Foundations and Trends in Signal Processing}, vol.~11, no.~1--2,
pp.~1--153, 2017.

\bibitem{detone2018}
D.~DeTone, T.~Malisiewicz, and A.~Rabinovich,
``SuperPoint: Self-Supervised Interest Point Detection and Description,''
in \textit{Proc. CVPR Workshops}, 2018.

\bibitem{lindenberger2023}
P.~Lindenberger, P.-E.~Sarlin, and M.~Pollefeys,
``LightGlue: Local Feature Matching at Light Speed,''
in \textit{Proc. ICCV}, 2023.

\bibitem{barath2020}
D.~Barath, J.~Noskova, M.~Ivashechkin, and J.~Matas,
``MAGSAC++, a Fast, Reliable and Accurate Robust Estimator,''
in \textit{Proc. CVPR}, 2020.

\end{thebibliography}

\end{document}
